
This appendix collects the standing assumptions and the formal statements of
every result the body refers to, and nothing else. Identifiers (L1, T2, T9,
and so on) are the register identifiers used throughout the supplementary
repository, so each statement here can be matched to its derivation
document, its proof, and its numerical checks. The register itself, with
per-result proof status, is \texttt{mlxor-derivations/THEOREMS.md}.

\subsection{Standing assumptions}

\textbf{A1} (local quadratic scope): $\Phi$ is $C^3$ and strongly concave on
the operating interval, with expansions taken at $\hstar$ and third-order
remainder bounded by $L_3 = \sup|\Phi'''|$.
\textbf{A2} (exploration class): all policies are non-anticipating; A2-sym
adds $\E[u_t \mid \mathcal{F}_t] = 0$, A2-adaptive is
unrestricted, and costs are measured from the certainty-equivalent anchor.
\textbf{A3} (noise exogeneity): the observation noise $\zeta_t$ is i.i.d.,
mean zero, variance $\sigma^2$, independent of $\mathcal{F}_t$. A3$'$ relaxes
this to a feedback channel with gain $\phi$.
\textbf{A4} (trust region): $|\dev_t| \le r$ pathwise.
\textbf{A5} (stationary scale): $S_T = \Theta(T)$.
\textbf{A6} (stable base loop): the retraining modulus satisfies $m < 1$.

The local dynamics are $\dev_{t+1} = -m\,\dev_t + u_t$ with
$\dev_t = h_t - \hstar$, observations
$\tau_t = \tau(\hstar) - \eps\,\dev_t + \zeta_t$, and $\Phi$ satisfying
$\Phi'(\hstar) = \delta_0$ and $-\Phi'' = \gpo$ at the operating point.
$\Sxx$ denotes the centered design energy $\sum_{t\le T}(\dev_t -
\bar\dev)^2$.

\subsection{The exchange rate and its floors}

\begin{lemma}[Cost equivalence; L1]
\label{lem:cost}
Under A1--A3 with conditionally symmetric exploration, the expected
incremental performative cost satisfies
$C_T = \tfrac{1}{2}\gpo\,\E\!\big[\sum_{t\le T}\dev_t^2\big] + R_T$ with
$|R_T| \le \tfrac{L_3}{6}\,\E\sum_t |\dev_t|^3$, and the realized cost's
variance decomposes exactly as
$\Var(C_T^{\mathrm{real}}) = \delta_0^2\,\Var\!\big(\sum_t \dev_t\big)
+ \tfrac{\gpo^2}{4}\,\Var\!\big(\sum_t \dev_t^2\big)$
with vanishing cross term.
\end{lemma}

\begin{theorem}[Information saturation; T1]
\label{thm:sat}
With no deliberate exploration and $\rho(M)<1$, the total design energy of
the retraining trajectory equals $\dev_0^\top P\,\dev_0$, where $P$ solves
the discrete Lyapunov equation $P = I + M^\top P M$. Consequently the Fisher
information for the response Jacobian is bounded in every direction,
uniformly in the horizon. In the scalar case the cap is
$\dev_0^2/(1-m^2)$, strictly increasing in $m$: a fast-contracting loop
reveals least about its own performativity (C1.1). Under the noisy feedback
channel A3$'$ the trajectory additionally carries a stationary excitation
floor of $(\sigma/\gamma)^2/(1-m^2)$ per step.
\end{theorem}

The body's Theorem \ref{thm:exchange} is the asymptotic form of the following
pathwise identity, which is what is proved in Appendix \ref{app:proofs}.

\begin{theorem}[The exchange rate, pathwise; T2]
\label{thm:t2exact}
Under A1, A2-sym, A3 and A6, in the stationary regime, on every realization,
\[
\Var(\hat\eps \mid \mathrm{design}) \times \tfrac{1}{2}\gpo \sum_{t\le T}
\dev_t^2 \;=\; \tfrac{1}{2}\,\gpo\,\sigma^2\Big(1 +
\tfrac{T\,\bar{\dev}^2}{\Sxx}\Big)
\;=\; \tfrac{1}{2}\,\gpo\,\sigma^2\,\big(1 + O_P(1/T)\big),
\]
independent of the exploration intensity, the horizon, and the modulus.
\end{theorem}

\begin{theorem}[Minimax lower bound, deviation budget; T3]
\label{thm:t3}
Over all non-anticipating exploration policies with
$\sum_{t\le T} \dev_t^2 \le D$ almost surely and all estimators, biased
included, the minimax risk over an interval of width $w$ satisfies
$\inf \sup_{\eps} \E[(\hat\eps - \eps)^2] \ge \sigma^2/(D + \sigma^2
(\pi/w)^2)$, so the exchange-rate product is a floor, attained with equality
by symmetric probing with least squares.
\end{theorem}

\begin{lemma}[Exploitation and information; L2]
\label{lem:l2}
For the symmetric two-point response family with half-width $\delta$ and any
non-anticipating policy under the certainty-equivalent anchor,
$\E[\mathrm{gain}_T] \le g_1\,\delta^{3/2}\,T^{1/2}\,(\E S_T)^{3/4}
\sigma^{-1/2}$ with $g_1 = \beta\,|\hstar - \psi|$, via the identity
$\E|\E[s\mid\mathcal{F}_t]| = \TV(P^+_t, P^-_t)$ and Pinsker's inequality.
At the minimax scale this is an $O(T^{-1/2})$ fraction of the exploration
cost.
\end{lemma}

\begin{theorem}[Minimax lower bound, value budget; T4]
\label{thm:t4}
Under stationary-scale exploration and the certainty-equivalent anchor, for
every non-anticipating policy and every estimator,
$\sup \Var(\hat\eps)\times C_T \ge \gpo\,\sigma^2/27\,(1 - O(T^{-1/2}))$,
the supremum over the two-point response family at the minimax scale, the
constant arising from a Le Cam reduction optimized at $c = 2/3$. The sharp
constant $\tfrac{1}{2}$, against which no simulated policy fell by more than
Monte Carlo error, is conjectured and is part of the open problem OPEN-1. It
is not proved here or anywhere in the repository.
\end{theorem}

\subsection{Design and certification}

\begin{theorem}[A-optimal exploration; T5a]
\label{thm:t5a}
Under the value budget $\tr(\Gpo M) \le B$ on the exploration second moment
$M$, the A-optimal design is $M^{*} = (B/\tr\,\Gpo^{1/2})\,\Gpo^{-1/2}$,
with value $(\tr\,\Gpo^{1/2})^2/B$: explore where the objective is flat.
\end{theorem}

\begin{corollary}[Price of isotropic jitter; C5.1]
\label{cor:c51}
Isotropic exploration at matched budget overpays the A-optimal design by
exactly the curvature-dispersion factor
$F = d\,\tr(\Gpo)/(\tr\,\Gpo^{1/2})^2 \ge 1$, with equality if and only if
all curvature directions are equal.
\end{corollary}

\begin{theorem}[Chebyshev unidentifiability; T6]
\label{thm:t6}
The sensitivity system of the structural family
$\tau(h) = C_0 + C_1 e^{-ch}$ is an extended Chebyshev system. Any design
with fewer than three support points has singular Fisher information, and
since noiseless retraining trajectories have at most two effective support
points, no retraining trajectory at any modulus identifies the response
curvature $c$.
\end{theorem}

\begin{theorem}[Anchoring crossover; T7]
\label{thm:t7}
Under the value budget $B$ and horizon $T$, the budget-and-horizon-optimal
nonparametric secant estimator has
$\mathrm{MSE}_{\mathrm{np}} = \gpo\sigma^2/(2B) + (\tau'''(\hstar)/6)^2
(2B/(\gpo T))^2$, so anchoring to a structural family with misspecification
$\delta_{\mathrm{mis}}$ wins, to leading order, if and only if
$\delta_{\mathrm{mis}} < |\tau'''(\hstar)|\,B/(3\,\gpo\,T)$.
\end{theorem}

\begin{lemma}[Perturbed modulus; L4]
\label{lem:l4}
If the corrected update uses an estimated response with $C^1$ error $e$, the
closed-loop modulus satisfies $|\hat\rho - \rho| \le \eta\,\beta\,
(\|e\|_\infty + |\hstar-\psi|\,\|e'\|_\infty)$.
\end{lemma}

\begin{proposition}[Safety under pessimism; P7.1]
\label{prop:p71}
Given an anytime-valid confidence sequence for the response parameters at
level $1-\delta$, the pessimism rule, which keeps the worst-case modulus over
the confidence set below $1 - c_{\mathrm{margin}}$ and freezes the correction
otherwise, keeps the realized modulus within the margin for all $t$
simultaneously with probability at least $1-\delta$. The result is
conditional on the validity of the confidence sequence, and the register
records it as such.
\end{proposition}

\subsection{Structure-proofness}

\begin{theorem}[Structure-proofness of the floor, exact reach; T9]
\label{thm:t9full}
Fix the structural family $\tau_\theta(h) = C_0 + C_1 e^{-ch}$
($C_1 > 0$, $c > 0$), the anchor $\hstar > 0$, the sensitivity
$s(h) = \partial_\theta \tau_\theta(h)$, and the estimand gradient
$g = \partial_\theta\,\eps(\hstar) = -s'(\hstar)$. For a finitely supported
design $\mu$ define the normalized Cram\'er--Rao product
\[
R(\mu) \;=\; \big(g^\top M(\mu)^{-1} g\big)\!\int (h - \hstar)^2\,
d\mu, \qquad M(\mu) = \int s\,s^\top d\mu,
\]
so that $R(\mu) \ge 1$ states that the family-knowing Cram\'er--Rao product
$\Var(\hat\eps)\times C_T$ is at least the exchange-rate floor
$\tfrac{1}{2}\gpo\sigma^2$ under the local-quadratic (A1) cost, for unbiased
estimators and, in the local-asymptotic-minimax sense, for all regular
estimators.
(i) \emph{Within reach the floor is structure-proof:} if
$\mathrm{supp}(\mu) \subset [\hstar - 2/c,\,\infty)$, then either
$\eps(\hstar)$ is not estimable under $\mu$, in which case the bound is
infinite, or $R(\mu) \ge 1$, strictly whenever $\mu$ places mass off the two
equality points $\{\hstar,\ \hstar - 2/c\}$; and $\inf_\mu R(\mu) = 1$,
approached but not attained by symmetric three-point designs collapsing at
$\hstar$. In particular the floor binds on every trust-region design with
$r \le 2/c$.
(ii) \emph{The reach is exact:} the pointwise domination underlying (i) fails
strictly below $\hstar - 2/c$, and an explicit four-point design with a
single below-reach probe of weight $3\times 10^{-4}$ attains $R = 0.8517$,
computer-verified at 50-digit precision as check V9.2.
\end{theorem}
